% Options for packages loaded elsewhere
\PassOptionsToPackage{unicode}{hyperref}
\PassOptionsToPackage{hyphens}{url}
%
\documentclass[
  11pt,
]{article}
\usepackage{amsmath,amssymb}
\usepackage{iftex}
\ifPDFTeX
  \usepackage[T1]{fontenc}
  \usepackage[utf8]{inputenc}
  \usepackage{textcomp} % provide euro and other symbols
\else % if luatex or xetex
  \usepackage{unicode-math} % this also loads fontspec
  \defaultfontfeatures{Scale=MatchLowercase}
  \defaultfontfeatures[\rmfamily]{Ligatures=TeX,Scale=1}
\fi
\usepackage{lmodern}
\ifPDFTeX\else
  % xetex/luatex font selection
\fi
% Use upquote if available, for straight quotes in verbatim environments
\IfFileExists{upquote.sty}{\usepackage{upquote}}{}
\IfFileExists{microtype.sty}{% use microtype if available
  \usepackage[]{microtype}
  \UseMicrotypeSet[protrusion]{basicmath} % disable protrusion for tt fonts
}{}
\makeatletter
\@ifundefined{KOMAClassName}{% if non-KOMA class
  \IfFileExists{parskip.sty}{%
    \usepackage{parskip}
  }{% else
    \setlength{\parindent}{0pt}
    \setlength{\parskip}{6pt plus 2pt minus 1pt}}
}{% if KOMA class
  \KOMAoptions{parskip=half}}
\makeatother
\usepackage{xcolor}
\usepackage[margin=2.2cm]{geometry}
\setlength{\emergencystretch}{3em} % prevent overfull lines
\providecommand{\tightlist}{%
  \setlength{\itemsep}{0pt}\setlength{\parskip}{0pt}}
\setcounter{secnumdepth}{-\maxdimen} % remove section numbering
\ifLuaTeX
  \usepackage{selnolig}  % disable illegal ligatures
\fi
\IfFileExists{bookmark.sty}{\usepackage{bookmark}}{\usepackage{hyperref}}
\IfFileExists{xurl.sty}{\usepackage{xurl}}{} % add URL line breaks if available
\urlstyle{same}
\hypersetup{
  pdftitle={Worksheet 1: Fractions},
  hidelinks,
  pdfcreator={LaTeX via pandoc}}

\title{Worksheet 1: Fractions}
\author{}
\date{}

\begin{document}
\maketitle

\hypertarget{part-a-numerical-fractions-review-10-problems}{%
\subsection{Part A: Numerical Fractions Review (10
problems)}\label{part-a-numerical-fractions-review-10-problems}}

Simplify or evaluate each expression. Show your work.

\textbf{1.} \(\dfrac{3}{4} + \dfrac{5}{6}\)

\textbf{2.} \(\dfrac{7}{8} - \dfrac{2}{3}\)

\textbf{3.} \(\dfrac{5}{9} \times \dfrac{3}{10}\)

\textbf{4.} \(\dfrac{4}{5} \div \dfrac{2}{15}\)

\textbf{5.} \(\dfrac{2}{3} + \dfrac{1}{4} - \dfrac{1}{6}\)

\textbf{6.} \(\dfrac{7}{12} \times \dfrac{8}{21} \times \dfrac{3}{2}\)

\textbf{7.} \(\dfrac{\frac{3}{4}}{\frac{9}{16}}\) (complex fraction)

\textbf{8.} \(\dfrac{\frac{2}{5} + \frac{1}{3}}{\frac{1}{2}}\)

\textbf{9.} \(2\dfrac{1}{3} \times 1\dfrac{1}{2}\) (mixed numbers)

\textbf{10.} \(4\dfrac{2}{5} \div 1\dfrac{1}{10}\)

\begin{center}\rule{0.5\linewidth}{0.5pt}\end{center}

\hypertarget{part-b-algebraic-fractions---simplification-8-problems}{%
\subsection{Part B: Algebraic Fractions - Simplification (8
problems)}\label{part-b-algebraic-fractions---simplification-8-problems}}

Simplify each expression completely. Assume no denominators equal zero.

\textbf{11.} \(\dfrac{6x^2}{9x}\)

\textbf{12.} \(\dfrac{15a^3b^2}{25ab^4}\)

\textbf{13.} \(\dfrac{x^2 - 4}{x + 2}\)

\textbf{14.} \(\dfrac{x^2 - 9}{x^2 + 6x + 9}\)

\textbf{15.} \(\dfrac{2x^2 + 4x}{x^2 + 2x}\)

\textbf{16.} \(\dfrac{x^2 - 5x + 6}{x^2 - 4}\)

\textbf{17.} \(\dfrac{6x^2 - x - 2}{2x^2 + 3x - 2}\)

\textbf{18.} \(\dfrac{x^3 - 8}{x^2 - 4}\)

\begin{center}\rule{0.5\linewidth}{0.5pt}\end{center}

\hypertarget{part-c-algebraic-fractions---operations-8-problems}{%
\subsection{Part C: Algebraic Fractions - Operations (8
problems)}\label{part-c-algebraic-fractions---operations-8-problems}}

Perform the indicated operations and simplify completely.

\textbf{19.} \(\dfrac{2}{x} + \dfrac{3}{x}\)

\textbf{20.} \(\dfrac{5}{x+1} - \dfrac{2}{x+1}\)

\textbf{21.} \(\dfrac{3}{x} + \dfrac{2}{y}\)

\textbf{22.} \(\dfrac{x}{x+2} + \dfrac{2}{x-2}\)

\textbf{23.} \(\dfrac{1}{x-1} - \dfrac{1}{x+1}\)

\textbf{24.} \(\dfrac{x}{x^2-1} + \dfrac{1}{x+1}\)

\textbf{25.} \(\dfrac{2x}{x+3} \times \dfrac{x^2 - 9}{4x^2}\)

\textbf{26.} \(\dfrac{x^2 - 4}{x + 5} \div \dfrac{x - 2}{x^2 + 5x}\)

\begin{center}\rule{0.5\linewidth}{0.5pt}\end{center}

\hypertarget{part-d-complex-algebraic-fractions-4-problems}{%
\subsection{Part D: Complex Algebraic Fractions (4
problems)}\label{part-d-complex-algebraic-fractions-4-problems}}

Simplify each complex fraction.

\textbf{27.}
\(\dfrac{\frac{1}{x} + \frac{1}{y}}{\frac{1}{x} - \frac{1}{y}}\)

\textbf{28.} \(\dfrac{\frac{2}{x+1} - \frac{1}{x}}{\frac{3}{x}}\)

\textbf{29.} \(\dfrac{1 - \frac{1}{x^2}}{1 + \frac{1}{x}}\)

\textbf{30.} \(\dfrac{\frac{a}{b} + \frac{b}{a}}{\frac{a^2 - b^2}{ab}}\)

\begin{center}\rule{0.5\linewidth}{0.5pt}\end{center}

\hypertarget{part-e-reasoning-about-fraction-behavior-6-problems}{%
\subsection{Part E: Reasoning About Fraction Behavior (6
problems)}\label{part-e-reasoning-about-fraction-behavior-6-problems}}

Answer each question with a brief explanation. No calculation required.

\textbf{31.} Consider the fraction \(\dfrac{n}{10}\) where \(n\) is a
positive number. - (a) What happens to the value of the fraction as
\(n\) gets larger and larger? - (b) What value does the fraction
approach as \(n\) approaches 0?

\textbf{32.} Consider the fraction \(\dfrac{10}{n}\) where \(n\) is a
positive number. - (a) What happens to the value as \(n\) gets larger
and larger? - (b) What happens to the value as \(n\) gets closer and
closer to 0 (but stays positive)?

\textbf{33.} For the fraction \(\dfrac{x + 1}{x}\) where \(x > 0\): -
(a) Rewrite this as a sum of two simpler fractions. - (b) As \(x\) gets
very large (like 100, 1000, 10000\ldots), what value does this fraction
approach? - (c) Why does this happen?

\textbf{34.} Consider \(\dfrac{x^2}{x^2 + 1}\) for positive values of
\(x\). - (a) Calculate the value when \(x = 1\), \(x = 10\), and
\(x = 100\). - (b) What value does this fraction seem to approach as
\(x\) grows without bound? - (c) Will the fraction ever actually reach
that value? Why or why not?

\textbf{35.} You have the fraction \(\dfrac{a}{b}\) where both \(a\) and
\(b\) are positive. - (a) If you double both the numerator and
denominator, what happens to the value? - (b) If you add 1 to both the
numerator and denominator, does the value increase, decrease, or stay
the same? (Hint: test with \(\frac{1}{2}\) and \(\frac{3}{4}\)) - (c)
Can you explain the pattern from part (b)?

\textbf{36.} The expression \(\dfrac{1}{1 + \frac{1}{x}}\) can be
simplified. - (a) Simplify this expression to a single fraction. - (b)
As \(x\) approaches infinity, what does the simplified expression
approach? - (c) As \(x\) approaches 0 (from the positive side), what
happens to the value?

\begin{center}\rule{0.5\linewidth}{0.5pt}\end{center}

\hypertarget{part-f-multi-step-problems-4-problems}{%
\subsection{Part F: Multi-Step Problems (4
problems)}\label{part-f-multi-step-problems-4-problems}}

These problems require multiple operations. Show all steps.

\textbf{37.} Simplify completely:
\(\dfrac{x^2 - 1}{x^2 + 2x} \cdot \dfrac{x^2 + 4x + 4}{x^2 + x - 2} \div \dfrac{x + 2}{x}\)

\textbf{38.} Express as a single fraction in lowest terms:
\[\dfrac{1}{x} + \dfrac{1}{x+1} + \dfrac{1}{x+2}\]

\textbf{39.} If \(\dfrac{1}{a} + \dfrac{1}{b} = \dfrac{1}{c}\), solve
for \(c\) in terms of \(a\) and \(b\).

\textbf{40.} The sum of a number and its reciprocal is \(\dfrac{5}{2}\).
- (a) Write an equation representing this situation. - (b) Multiply
through to clear fractions. - (c) Find the number(s).

\begin{center}\rule{0.5\linewidth}{0.5pt}\end{center}

\end{document}
